\documentclass[11pt, aspectratio=169]{beamer}
\usepackage[style=authoryear-icomp, backend=biber]{biblatex}
\usepackage{booktabs}
\usetheme{jcar}

%\includeonlyframes{a} % Para compilar solo en la que se está trabajando

%\addbibresource{refs.bib}

\begin{document}
\author[Juan C. Arboleda Rivera]{Juan Camilo Arboleda Rivera}
\title{Ejemplo práctico:\\Comparación de dos muestras}
%\subtitle{}
%\titlegraphic{\includegraphics[width=2cm]{udea}}
\institute[U. de A.]{Instituto de Biología\\
  Facultad de Ciencias Exactas y Naturales\\
  Universidad de Antioquia\\
  juan.arboleda2@udea.edu.co}
\date{v: \today}
%\subject{}

\begin{frame}[plain]
  \maketitle
\end{frame}

%\begin{frame}
% \frametitle{Contenidos}
% \tableofcontents
%\end{frame}

\section{}

% ------------------------------------------------
\begin{frame}{Una pregunta biológica}

Supongamos que queremos saber si un antibiótico afecta el crecimiento de una
bacteria.

\vspace{0.5cm}

Tenemos dos grupos:

\begin{itemize}
    \item \textbf{Control}: bacterias sin antibiótico.
    \item \textbf{Tratamiento}: bacterias con antibiótico.
\end{itemize}

\end{frame}

% -------------------------------------------------
\begin{frame}
[label=a]
{Densidad óptica (OD)}{}
 \centering
 Medimos el crecimiento de cada grupo mediante la densidad óptica a 600 nm:

 \vspace{-1em}

 \[
 \text{OD}_{600}
 \]

 \includegraphics[width=\linewidth]{absorption_vs_od.png}
\end{frame}

% -------------------------------------------------
\begin{frame}
[label=a]
{Densidad óptica (OD)}{}
 \centering
 \includegraphics[height=.85\textheight]{densidad-óptica.jpeg}
\end{frame}

% -------------------------------------------------
\begin{frame}
[label=a]
{}{}
 \centering
 \includegraphics[height=.85\textheight]{EcoliSaturated-12.png}
\end{frame}

\begin{frame}
[label=a]
{}{}
 \centering


 \begin{alertblock}{Pregunta}
 ¿La diferencia que observamos entre los grupos es suficientemente grande
 como para pensar que el tratamiento realmente tiene un efecto?
 \end{alertblock}

% \includegraphics[height=.85\textheight]{}
\end{frame}

% ------------------------------------------------
\begin{frame}{Los datos}

\begin{columns}
 \column{.5\linewidth}
  \begin{table}
  \centering
  \begin{tabular}{cc}
  \toprule
  Control & Tratamiento \\
  \midrule
  0.82 & 0.61\\
  0.76 & 0.58\\
  0.89 & 0.69\\
  0.64 & 0.73\\
  0.88 & 0.56\\
  0.79 & 0.62\\
  0.84 & 0.59\\
  0.67 & 0.65\\
  0.87 & 0.57\\
  0.80 & 0.60\\
  \bottomrule
  \end{tabular}
  \end{table}
 \column{.5\linewidth}
  Las medias son aproximadamente:

  \[
  \bar{x}_{C}=0.796
  \qquad
  \bar{x}_{T}=0.620
  \]

  Por tanto,

  \[
  \boxed{
  \bar{x}_{C}-\bar{x}_{T}=0.176
  }
  \]
\end{columns}

\end{frame}

% ------------------------------------------------
\begin{frame}{Pero... ¿0.176 es mucho o poco?}

Una diferencia entre dos grupos puede aparecer por diferentes razones.

\vspace{0.3cm}

\begin{enumerate}[<+->]
    \item Existe variabilidad biológica entre las muestras.
    \item La diferencia apareció por azar.
    \item El tratamiento realmente tiene un efecto.
\end{enumerate}

\vspace{0.5cm}

\onslide<4->{

 La estadística nos ayuda a responder:

 \begin{alertblock}{Pregunta clave}
 Si el tratamiento \textbf{no tuviera ningún efecto},
 ¿qué tan frecuente sería observar una diferencia como la que encontramos?
 \end{alertblock}
}

\end{frame}

% ------------------------------------------------
\begin{frame}{La hipótesis nula}

Para formalizar la pregunta planteamos una hipótesis nula:

\only<1>{\centering ¿Cuál sería la hipótesis en este caso?}

\pause

\[
H_0:
\quad
\text{El tratamiento no cambia el crecimiento bacteriano}
\]

Es decir, bajo $H_0$:

\[
\boxed{\text{Control y tratamiento son intercambiables}}
\]

\vspace{0.5cm}

Si esto fuera cierto, ¿qué impediría que una bacteria que estaba
en el grupo control estuviera en el grupo tratamiento?

\vspace{0.3cm}

\pause

Nada.

\begin{center}
\Large
Podemos \textbf{reordenar aleatoriamente} las etiquetas.
\end{center}

\end{frame}

% ------------------------------------------------
\begin{frame}{La idea de la permutación}

Tenemos las 20 mediciones:

\[
0.82,\;0.76,\;0.89,\ldots,\;0.57,\;0.60
\]

Ignoramos temporalmente las etiquetas
``control'' y ``tratamiento''.

\vspace{0.4cm}

Después:

\begin{enumerate}
    \item Mezclamos las 20 observaciones.
    \item Elegimos 10 al azar para formar un nuevo grupo.
    \item Las otras 10 forman el segundo grupo.
    \item Calculamos nuevamente la diferencia entre las medias.
\end{enumerate}

\vspace{0.3cm}

Repetimos este procedimiento muchas veces.

\end{frame}

% ------------------------------------------------
\begin{frame}{¿Qué estamos simulando?}

Cada reorganización responde a la pregunta:

\begin{center}
\Large
``¿Qué diferencia podríamos observar \\
si $H_0$ fuera cierta?''
\end{center}

\vspace{0.5cm}

Después de muchas permutaciones obtenemos una distribución:

\[
\boxed{
\text{Distribución de diferencias esperadas bajo }H_0
}
\]

Esta distribución representa lo que podría ocurrir
\textbf{por azar} si el tratamiento no tuviera efecto.

\end{frame}

% ------------------------------------------------
\begin{frame}{La diferencia observada}

Nuestra diferencia original fue:

\[
D_{\mathrm{obs}}
=
\bar{x}_{C}-\bar{x}_{T}
=
0.176
\]

Ahora imaginemos que hacemos:

\[
10\,000
\]

permutaciones.

\vspace{0.4cm}

Para cada permutación calculamos:

\[
D_i
=
\bar{x}_{1,i}-\bar{x}_{2,i}
\]

y preguntamos:

\begin{center}
\Large
¿Cuántas diferencias son tan extremas como $0.176$?
\end{center}

\end{frame}

% ------------------------------------------------
\begin{frame}{¿Qué es el valor $p$?}

El valor $p$ es la proporción de resultados tan extremos
como el observado, \textbf{suponiendo que $H_0$ es cierta}.

\vspace{0.5cm}

En una prueba de permutación podemos estimarlo mediante:

\[
\boxed{
\hat p =
\frac{\#\{\text{permutaciones tan extremas como la observada}\}}
{\#\{\text{permutaciones}\}}
}
\]

\end{frame}

% ------------------------------------------------
\begin{frame}{Un ejemplo}

Supongamos que hacemos:

\[
10\,000
\]

permutaciones.

Encontramos que solamente: \vspace{-1em}

\[
73
\]

producen una diferencia igual o mayor que la observada
en valor absoluto.

Entonces:

\[
\hat p =
\frac{73}{10\,000}
=
0.0073
\]

Por tanto:

\[
\boxed{p \approx 0.007}
\]

\end{frame}

% ------------------------------------------------
\begin{frame}{¿Cómo interpretamos $p=0.007$?}

Una interpretación adecuada sería:

\vspace{0.4cm}

\begin{alertblock}{Interpretación}
Si realmente no hubiera diferencia entre control y tratamiento,
una diferencia tan extrema como la observada aparecería en
aproximadamente el $0.7\%$ de las reorganizaciones aleatorias.
\end{alertblock}

\vspace{0.5cm}

Esto proporciona evidencia contra la hipótesis nula.

\end{frame}

% ------------------------------------------------
\begin{frame}{Lo que el valor $p$ NO significa}

Si obtenemos:

\[
p=0.007
\]

\textbf{NO} significa:

\begin{itemize}
    \item ``Hay un $0.7\%$ de probabilidad de que $H_0$ sea cierta.''
    \item ``Hay un $99.3\%$ de probabilidad de que el tratamiento funcione.''
    \item ``El resultado es $99.3\%$ correcto.''
\end{itemize}

\vspace{0.4cm}

El valor $p$ habla de:

\begin{center}
\Large
\textbf{los datos bajo la hipótesis nula}
\end{center}

no de la probabilidad de que la hipótesis sea verdadera.

\end{frame}

% ------------------------------------------------
\begin{frame}{¿Por qué usamos un umbral como 0.05?}

Es común utilizar:

\vspace{-1em}

\[
\alpha=0.05
\]

como umbral de decisión.

\vspace{0.5cm}

Una regla convencional es:

\[
\begin{cases}
p < 0.05 &\Rightarrow \text{rechazar }H_0\\
p \geq 0.05 &\Rightarrow \text{no rechazar }H_0
\end{cases}
\]

Pero:

\begin{alertblock}{Importante}
$0.05$ no es una frontera mágica entre
``verdadero'' y ``falso''.
\end{alertblock}

\end{frame}

% ------------------------------------------------
\begin{frame}{Significativo no significa importante}

Imaginemos dos experimentos.

\vspace{0.3cm}

\textbf{Experimento A}

\[
\Delta = 0.20,
\qquad p=0.001
\]

\textbf{Experimento B}

\[
\Delta = 0.02,
\qquad p=0.001
\]

\vspace{0.4cm}

Ambos pueden ser estadísticamente significativos.

Pero la pregunta biológica es diferente:

\begin{center}
\large
¿Qué tan grande es el efecto y qué significa biológicamente?
\end{center}

\end{frame}

% ------------------------------------------------
\begin{frame}{Valor $p$ y tamaño del efecto}

Una buena pregunta científica no es solamente:

\[
\text{``¿Es significativo?''}
\]

También debemos preguntar:

\[
\boxed{\text{¿Qué tan grande es la diferencia?}}
\]

y

\[
\boxed{\text{¿Es biológicamente relevante?}}
\]

\vspace{0.5cm}

En nuestro ejemplo:

\vspace{-1em}

\[
\Delta = 0.176
\]

El valor $p$ nos ayuda a evaluar la evidencia contra $H_0$,
pero no reemplaza la descripción del efecto.

\end{frame}

% ------------------------------------------------
\begin{frame}{Ejercicio}

Usen los datos del experimento bacteriano.

\vspace{0.3cm}

\begin{enumerate}
    \item Mezclen las 20 observaciones.
    \item Asignen aleatoriamente 10 al grupo A y 10 al grupo B.
    \item Calculen:
    \[
    D=\bar{x}_A-\bar{x}_B
    \]
    \item Repitan el procedimiento muchas veces.
    \item Construyan un histograma de las diferencias.
    \item Marquen la diferencia observada:
    \[
    D_{\mathrm{obs}}=0.176
    \]
    \item Estimen el valor $p$.
\end{enumerate}

\end{frame}

% ------------------------------------------------
\begin{frame}{Resumen}

\begin{enumerate}
    \item El valor $p$ se calcula \textbf{suponiendo que la hipótesis nula
    es cierta}.

    \item Pregunta fundamental:
    \[
    \text{¿Qué tan compatible son los datos con }H_0?
    \]

    \item Una prueba de permutación permite responder esta pregunta
    de manera intuitiva mediante la aleatorización.

    \item Un valor $p$ pequeño proporciona evidencia contra $H_0$,
    pero no mide el tamaño ni la importancia biológica del efecto.
\end{enumerate}

\vspace{0.4cm}

\begin{center}
\Large
\textbf{Datos $\rightarrow$ azar bajo $H_0$ $\rightarrow$ valor $p$
$\rightarrow$ interpretación biológica}
\end{center}

\end{frame}

%\appendix
%\AtNextBibliography{\footnotesize}
%\begin{frame}[plain, allowframebreaks]{Bibliografía}{}
% \bibitemsep 1em
% \printbibliography[heading=none]
%\end{frame}
\end{document}
